\subsection{5. Allgemeine CAR-Identität auf
N=3}\label{allgemeine-car-identituxe4t-auf-n3}

Bezeichne die globale Dreifermion-zu-Boson-Fermion-Abbildung mit
\(T_C\). Die Boson-Fermion-Zeilen werden als \((x,A,i)\) angeordnet. Aus
der Paar-CAR-Kontraktion folgt

\[\adjustbox{max width=\linewidth}{$\displaystyle 
\boxed{G_C^{(3)}:=T_CT_C^\dagger
=8I_L\otimes I_{3840}-(C^\dagger C)\otimes K.}
$}\]

Für die gewählte Vernichtungskonvention ist der Sitefaktor \textbf{C†C},
nicht CC†. Aus dem Austauschterm folgt komponentenweise

\[\adjustbox{max width=\linewidth}{$\displaystyle 
\sum_z\overline{C_{xz}}C_{yz}
=(\overline C C^T)_{xy}=(C^\dagger C)_{xy},
$}\]

weil C symmetrisch ist. Der direkte Term ist wegen der
Frobeniusnormierung 8I. Der innere Austauschblock ist genau das oben
unabhängig geprüfte K. Das beweist die allgemeine Identität. Für reelles
C sind beide scheinbar verschiedenen Konventionen identisch; deshalb
wird zusätzlich ein komplexes Gegenbeispiel
\(C=\left(\begin{smallmatrix}1&i\\i&0\end{smallmatrix}\right)/\sqrt3\)
durch vollständigen CAR-Aufbau geprüft. Reale und imaginäre Teile werden
separat ganzzahlig multipliziert.

Sind γ die Eigenwerte von C†C und λ die von G\_0, so sind die Eigenwerte
des globalen Gramoperators

\[\adjustbox{max width=\linewidth}{$\displaystyle 
8-\gamma(8-\lambda),
$}\]

mit den entsprechenden Produktmultiplizitäten. Bei L=2 ergibt sich:

{\def\LTcaptype{none} % do not increment counter
\begin{longtable}[]{@{}ll@{}}
\toprule\noalign{}
Fall & Vollständiges Spektrum von \(G_C^{(3)}\) \\
\midrule\noalign{}
\endhead
\bottomrule\noalign{}
\endlastfoot
Verteilt & 4 {[}128{]}; 15/2 {[}5760{]}; 9 {[}1152{]}; 10 {[}640{]} \\
Kollektiv & 0 {[}64{]}; 7 {[}2880{]}; 8 {[}3840{]}; 10 {[}576{]}; 12
{[}320{]} \\
\end{longtable}
}

Der Prüfer baut beide Abbildungen auf allen \(\binom{128}{3}=341376\)
Dreifermionzuständen und 7680 Boson-Fermionzuständen direkt auf. Die
zugrundeliegenden vollständigen N=3-Hamiltonoperatoren haben jeweils
349056 Dimensionen. Die Rechnung braucht nur dünn besetzte
Gramoperatoren, keine große dichte Diagonalisierung.

Für Δ\textgreater0 und g≠0 ist der Eigenraum von H\_C bei der
\textbf{exakten} Energie Δ isomorph zu \(\ker G_C^{(3)}\). Denn die
Eigenwertgleichungen lauten für Paarteil u und Bosonteil v

\[\adjustbox{max width=\linewidth}{$\displaystyle 
gT_C^\dagger v=\Delta u,\qquad gT_Cu=0.
$}\]

Daraus folgen \(G_C^{(3)}v=0\), dann \(T_C^\dagger v=0\) und u=0. Also
hat das verteilte Modell \textbf{keinen} solchen Zustand, das kollektive
Modell \textbf{genau 64}. Beide gesamten N=3-Spektren sind verschieden.
Eine Umbenennung von Ports oder eine passive Basiswahl kann diesen
spektralen Unterschied nicht entfernen.

\subsection{6. Positive Rückrechnung: der Schatten bestimmt
Q}\label{positive-ruxfcckrechnung-der-schatten-bestimmt-q}

Die obige Identität lässt sich tatsächlich invertieren. Mit
\(Q=C^\dagger C\) gilt

\[\adjustbox{max width=\linewidth}{$\displaystyle 
\boxed{Q=\frac1{960}\operatorname{Tr}_{\mathbb C^{60}\otimes\mathbb C^{64}}
\left(8I-G_C^{(3)}\right).}
$}\]

Der Teiltrace ist \textbf{unnormiert} und lässt den Kopienfaktor stehen.
Die Zahl 960 ist die exakt aus W gerechnete Spur von K. Für beide
vollständigen L=2-Gegenmodelle wird Q eintragsweise aus dem direkt
aufgebauten Gramoperator wiedergewonnen.

Das ist die gesuchte Art von Brücke: Eine höhere dynamische Antwort
verrät etwas über die gemeinsame Zusammensetzung, das die gesamte
niedrigere Spektralantwort nicht verrät. Sie liefert Q innerhalb der
erklärten faktorisierten Familie. Sie wählt weder C aus dem Compiler aus
noch rekonstruiert sie alle phasenmarkierten Operationen oder eine
einzigartige Welt. Insbesondere kann dasselbe Q mehrere markierte
C-Matrizen zulassen. Nach vollständigem unmarkiertem unitärem
Kopienquotienten bestimmen die Singulärwerte die Takagi-Klasse; ein
solcher Quotient muss zum Operationsvertrag passen.

\subsubsection{Anschluss an dieselbe
Bosonzahl-Auslese}\label{anschluss-an-dieselbe-bosonzahl-auslese}

Auf N=3 hat H\_C die Blockform
\(\left(\begin{smallmatrix}0&gT_C^\dagger\\gT_C&\Delta I\end{smallmatrix}\right)\).
Für einen normierten Eingang v auf der Boson-Fermion-Seite gilt exakt im
Taylor-Anfang

\[\adjustbox{max width=\linewidth}{$\displaystyle 
p_{N_b=1}(t)=1-g^2\langle v,G_C^{(3)}v\rangle t^2+O(t^4).
$}\]

Die erste und dritte Ableitung verschwinden. Der Protokollstrang hat
diesen Blocksatz unabhängig für beliebige endliche Blockdimensionen
symbolisch geprüft. Durch zugängliche Superpositionen und Polarisation
könnte man die Matrixelemente aus der Anfangskrümmung rekonstruieren und
dann Q berechnen. \textbf{Vorausgesetzt} wären solche Präparationen,
eine kalibrierte Bosonzahlmessung, g und ein Zeitmaß. Bloße
Spektraleigenwerte liefern nur die γ-Werte, nicht das markierte Q.

Dies ist eine echte gemeinsame Operatorverbindung zwischen
Kompositionsrechnung und Aufzeichnungsrechnung. Die konkreten Abläufe
liegen trotzdem in verschiedenen Ladungssektoren und sind noch kein
gemeinsamer nativer Grundzustandsversuch.

\subsection{7. Warum ein einzelnes C möglicherweise zu einfach für
gerichteten Ladungstransport
ist}\label{warum-ein-einzelnes-c-muxf6glicherweise-zu-einfach-fuxfcr-gerichteten-ladungstransport-ist}

Jede komplex-symmetrische endliche Matrix besitzt eine Takagi-Zerlegung
\(U^T C U=D\) mit unitärem U und nichtnegativem diagonalem D. Eine
Standardreferenz ist \href{https://arxiv.org/abs/2009.03990}{Haber, A
tale of three diagonalizations, Abschnitt 7}. Der Prüfer kontrolliert
zusätzlich ein nichtreelles rationales Beispiel. Der allgemeine
Existenzsatz wird aus der linearen Algebra übernommen, nicht aus diesem
Beispiel abgeleitet.

In der entsprechenden Fermionbasis enthält jeder Vertex zwei Fermionen
aus derselben Takagi-Kopie. Deshalb bleibt \textbf{jede dieser
Kopienparitäten} unter H\_C, X, Nb und innerem G erhalten. Das
gemeinsame Bosonsystem kann weiterhin koppeln und korrelieren; dieser
Satz bedeutet nicht vollständige Faktorisierung des gesamten Modells. Er
verhindert aber einen behaupteten einzelnen Fermionwechsel zwischen
genau diesen konservierten Kopiensektoren mit diesem Alphabet.

Die Takagi-Achsen sind nicht automatisch räumliche Orte. Für andere
markierte, nichtorthogonale Teile folgt kein universelles
Transferverbot. Mehrere nicht gleichzeitig diagonalisierbare
Kopplungstensoren können diese spezielle Paritätszerlegung aufbrechen.
Das Zwei-Chart-Modell liefert genau ein solches kleines Beispiel: Die
beiden Rang-eins-Projektoren auf e1 und (3/5,4/5) kommutieren nicht und
besitzen nur den skalaren gemeinsamen Kommutanten. Ihre Auswahl bleibt
zusätzliche Komposition.

\subsection{8. Präzise nächste
Entscheidung}\label{pruxe4zise-nuxe4chste-entscheidung}

Der einfachste nächste Auftrag ist kein neuer großer Tensorraum. Er ist
eine typisierte Quellvorschrift für \textbf{zwei} Einbettungen E\_L,E\_R
und ihre Überlappung \(E_L^\dagger E_R\), einschließlich der
Bosonbankzuordnung. Die vorhandenen Quellstrahlen enthalten diskrete
Überlappungszahlen, aber der direkte Import ihrer internen
C4-Projektoren als neue triviale Kopien ist nicht äquivariant begründet.

Wenn der nächste Quellschritt existiert, muss er die hier berechneten
Antworten vorhersagen, ohne c oder C anschließend anzupassen. Wenn er
nur symmetrieverträgliche freie C zulässt, ist die Herkunftsfrage nicht
geschlossen. Die Rückrechnungsformel stellt dann einen Mess- oder
Identifikationsvertrag bereit, keinen Ersatz für die Herleitung.

\subsection{9. Reproduktion und Grenze}\label{reproduktion-und-grenze}

\codeword{verify\_composition.py} führt 67 explizite exakt arbeitende
Prüfbedingungen aus, einschließlich Quellenpins und Strukturguards. Es
gibt keine numerischen Toleranzprüfungen. Die allgemeine Formel und ihre
Inversion sind oben algebraisch hergeleitet; drei vollständige
N=3-CAR-Aufbauten kontrollieren Normalisierung, Vorzeichen und komplexe
Konvention unabhängig.

Ein neuer Grundzustand für diese erweiterten Familien wurde nicht
berechnet. Die früheren nativen N=64- und Polschranken dürfen nicht auf
sie übertragen werden. Es ist keine physische Raumdimension, keine
lokale Kausalstruktur, kein chirales Maß und keine Gravitation
hergeleitet. Alle T1-T8-Gesamtpflichten bleiben offen.

\begin{center}\rule{0.5\linewidth}{0.5pt}\end{center}

\section{Teilbericht A: ursprüngliche Quelle und überlappende
Charts}\label{teilbericht-a-urspruxfcngliche-quelle-und-uxfcberlappende-charts}

\section{Quellenkomposition: dieselben lokalen Vertices, verschiedene
gemeinsame
Fermionen}\label{quellenkomposition-dieselben-lokalen-vertices-verschiedene-gemeinsame-fermionen}

\begin{enumerate}
\def\labelenumi{\arabic{enumi}.}
\setcounter{enumi}{14}
\tightlist
\item
  September 2026 · v1.6.9 · NON-RH · isolierter Quellenstrang
\end{enumerate}

\subsection{Stärkstes neues Ergebnis}\label{stuxe4rkstes-neues-ergebnis}

Die geprüfte Quelle bestimmt bereits \textbf{einen gemeinsamen lokalen
CAR-Träger mit 64 Fermionmoden}, nicht 480 unabhängige
Wechselwirkungsbanken. Jeder dieser Fermionmoden kommt in genau 15 der
480 nichtverschwindenden Vertices vor. Aber die Quelle bestimmt in den
untersuchten Konstruktionen \textbf{noch keine Überlappung zweier
vollständiger Kopplungs-Charts dieses Trägers}.

Das lässt sich durch ein genau geschlossenes Gegenmodell sichtbar
machen: Zwei Modelle auf derselben Fockbasis mit 128 Fermion- und 120
Bosonmoden behalten jeden lokalen W-Koeffizienten, dieselbe Kopplung g,
dieselbe Bosonenergie Δ, volles G = Spin(10) × SU(4), denselben passiven
Clock und dieselbe erhaltene Gesamtzahl. Beide benutzen wirklich
gemeinsame Fermionressourcen. Trotzdem ergeben sie verschiedene
Übergangswahrscheinlichkeiten. Eine dimensionslose Kombination der
ersten vier Energiemomente liest die zusätzlich gewählte Überlappung
exakt wieder aus.

\textbf{Status:} exakter Satz über die hier angegebene Familie; kein
Satz über alle Compiler, keine ausgewählte Raumgeometrie und keine
Herleitung physischer Präparation oder Zeit.

\subsection{1. Was die gelesenen Quellen tatsächlich
zusammensetzen}\label{was-die-gelesenen-quellen-tatsuxe4chlich-zusammensetzen}

Alle genannten Dateien sind in \href{inputs/}{inputs} eingefroren;
Originalpfade, Bytezahlen und SHA-256 stehen im
\href{inputs_manifest.json}{Manifest}. Der Prüfer importiert keine
beweglichen Repository-Dateien.

{\def\LTcaptype{none} % do not increment counter
\begin{longtable}[]{@{}
  >{\raggedright\arraybackslash}p{(\linewidth - 4\tabcolsep) * \real{0.3333}}
  >{\raggedright\arraybackslash}p{(\linewidth - 4\tabcolsep) * \real{0.3333}}
  >{\raggedright\arraybackslash}p{(\linewidth - 4\tabcolsep) * \real{0.3333}}@{}}
\toprule\noalign{}
\begin{minipage}[b]{\linewidth}\raggedright
Quellschritt
\end{minipage} & \begin{minipage}[b]{\linewidth}\raggedright
Tatsächlich gegeben
\end{minipage} & \begin{minipage}[b]{\linewidth}\raggedright
Nicht durch diesen Schritt gegeben
\end{minipage} \\
\midrule\noalign{}
\endhead
\bottomrule\noalign{}
\endlastfoot
\codeword{verification/v566\_parabolic\_anchor\_selfcode.py}, Konstanten
111-121 und S1 & U und V auf \textbf{demselben C³}; die sieben Wörter
I,U,V,UV,VU,V²,VUV spannen eine 7-dimensionale, unter allen 49 Produkten
geschlossene Algebra. Neu exakt nachgerechnet. & Neue Träger,
Fockfaktoren, räumliche Nachbarschaft oder eine CAR-Einbettung von C³ in
das spätere Modell. \\
\codeword{verification/v783\_two\_qubit\_clifford.py}, P0/P1;
\codeword{compiler-origin-audit-20260913/context\_instrument.py}, 69-114
& 60 konkrete Gaußwurzelstrahlen auf \textbf{demselben C⁴}, ihre
Projektoren und exakten Überlappungen. Neu aus dem gepinnten
Originalpräfix rekonstruiert. & Die Identifikation dieses C⁴ mit einer
G-trivialen Kopien-/Multiplizitätsebene des 64-Fermionmodells. \\
\codeword{universalraum-singlet-observable-20260915/sources/native\_source.py},
36-85, 132-220 & Aus fünf Hilfs-Jordan-Wigner-Moden, geradem Spinor und
antisymmetrischem Farbfaktor entsteht W: Λ²C⁶⁴ → C⁶⁰. Die
\textbf{physischen} 64 Fermionmoden werden gemeinsam benutzt. WWᵀ=8I;
alle 60 Lie-Intertwiner werden geprüft. & Ein Ortsindex x, eine Zahl von
Kopien, deren gegenseitige Antikommutatoren oder ein globaler
Site-Tensor. Die fünf Hilfs-CAR-Moden des Cliffordbaus sind nicht die 64
physischen f-Moden. \\
\codeword{universalraum-native-operations-ground-response-20260915/common.py},
1-17 & Das deklarierte native H und N=Nf+2Nb auf dieser einen Fockbank.
& Ein Eindeutigkeitssatz, dass die ursprüngliche Compilergrammatik genau
dieses physische H, g/Δ oder seinen Mehrteilträger erzwingt. \\
\codeword{compiler-origin-audit-20260913/record\_composition.py} &
Gegebene Kontext-Vormessungen lassen sich als Matrizen multiplizieren;
kohärent wiederverwendete und frische Register sind unterschiedliche
Prozesse. Schon der Header nennt die zusätzlichen Ressourcen. &
Lieferung frischer Register, inter-Register-Kopplung, Kontextpolitik,
Born-Regel oder physische Zeit. \\
\codeword{systematic-origin-audit-20260912/interaction/composition.py} &
Bedingt gesetzte Bell-Kanten teilen einen echten Faktor; überlappende
Projektoren sind nicht unabhängig. & Physischer Graph und
Tensorfaktorisierung. Der Header grenzt das selbst ausdrücklich ab. \\
\codeword{universalraum-fugen-20260914/checker.py};
\codeword{universalraum-v14-F2-architektur-20260914/README.md} und
\codeword{checker.py} & Explizite Modelle für interne Clebsch-Kanten,
Tetramerzellen, Zell- oder Kantenbanken. Die F2-Auswahl gilt für die
deklarierte harte C16-/Matching-/diagonale Gaugeklasse ohne Hopping. &
Eine unbedingte Auswahl derselben Architektur im v1.6.x-Voll-G-Modell.
Die README nennt ausdrücklich als fehlenden möglichen Quellschritt eine
Vorschrift, die K4-Zellen oder eine Kette erzwingt. \\
\end{longtable}
}

Die Quellen enthalten also echte Multiplikation, Überlappung und interne
Verklebung. Die fehlende Kante ist \textbf{nicht „es gibt überhaupt
keine Komposition``}, sondern der typisierte Übergang von diesen
festliegenden internen Trägern zu globalen gemeinsam benutzten
CAR-Ressourcen.

\subsubsection{Clock ist keine automatische
Vervielfältigung}\label{clock-ist-keine-automatische-vervielfuxe4ltigung}

Der ursprüngliche Clock wird aus
\codeword{experiments/tfpt-discovery/seam\_state\_derivation\_probe.py},
S0.1-S0.4, neu ausgeführt. Sein dokumentierter passiver Lift erfüllt mit
L=Λ²GF:

W L = GB W, Lᵏ Wᵀ = Wᵀ GBᵏ, k=0,\ldots,5.

Da GB orthogonal ist und W Rang 60 hat, sind alle sechs rotierten hellen
Räume genau derselbe 60-dimensionale Raum. Der Vereinigungsrang ist
\textbf{60, nicht 360}. Das wird durch sechs exakte Matrixidentitäten
bewiesen; der in \codeword{clock\_gluing.py} vorhandene
Mehrprimzahl-Rangheurismus wird dafür nicht verwendet. Interne
Clock-Wiederholung allein liefert deshalb keinen neuen unabhängigen
hellen Kopienträger.

\subsection{2. Kleinstes skalares
Überlappungs-Gegenmodell}\label{kleinstes-skalares-uxfcberlappungs-gegenmodell}

Seien a\_i,d\_i, i=0,\ldots,63, zwei kanonische Familien auf
\textbf{einer} gemeinsamen CAR-Fockbasis. Setze für
0\textless c\textless1 und s=√(1−c²):

f\_i\^{}L=a\_i, f\_i\^{}R=c a\_i+s d\_i.

Jeder Chart für sich ist kanonisch; gegenseitig gilt

\{f\_i\textsuperscript{L,f\_j}\{R†\}\}=c δij.

Das sind überlappende Charts, \textbf{keine unabhängigen räumlichen
Orte}. Ihre gemeinsame Gramform ist {[}{[}1,c{]},{[}c,1{]}{]}⊗I64 und
hat Rang 128. Für den erklärten skalaren, nichtdegenerierten
Überlappungsansatz sind 128 globale Fermionmoden daher minimal. Für
beide Vergleichsmodelle sind c und s ungleich null: Es wird kein
unbenutzter Fermion-Kopienraum als bloßer Zuschauer hinzugefügt.

Zusätzlich gibt es zwei unabhängige Bosonbanken b\_A,L und b\_A,R,
A=0,\ldots,59. Mit exakt demselben eingefrorenen W definieren wir

P\_A,x = Σ\_\{i\textless j\} W\_A,ij f\_j\^{}x f\_i\^{}x,

H\_c = Δ(Nb,L+Nb,R) + g Σ\_\{A,x\}(b\_A,x† P\_A,x + P\_A,x† b\_A,x).

Diese Summe ist die \textbf{deklarierte Gegenmodell-Kompositionsregel},
nicht ein behaupteter Originalcompiler-Output. Sie ist auf jedem festen
N-Sektor eine endliche hermitesche Matrix. Erhalten sind

N = Na+Nd+2(Nb,L+Nb,R), volles G, der diagonal wirkende ursprüngliche
Clock.

„Dieselben lokalen Daten`` bedeutet hier: Beide isolierten
Chart-Algebren und beide einzelnen Chart-Hamiltonoperatoren sind durch
kanonische Einbettungen exakt dieselben nativen Modelle, einschließlich
ihrer höheren lokalen Operatoridentitäten. Es bedeutet \textbf{nicht},
dass bei gleichzeitig eingeschalteter globaler Summe jeder bisherige
Einchart-Messwert unverändert bleibt. Gerade die gemeinsame
Ressourcenbenutzung ist zusätzliche Dynamik.

Die vollständige G-Kovarianz folgt aus allen 45+15 neu geprüften
W-Intertwinern und der exakten Identität (I₂⊗X)(u⊗I64)=(u⊗I64)X. Alle
global erweiterten Vertices werden zusätzlich gegen sämtliche acht
Cartangewichte geprüft. Es wird kein neues Zahl-veränderndes
Referenzsystem benötigt.

\textbf{Ressourcenbuch:} W und seine internen Symmetrien stammen aus dem
eingefrorenen nativen Konstrukt. Hinzugefügt sind zwei Chartnamen, ihre
Einbettungen und c, zwei Bosonbanknamen, die summierte H-Regel, sowie
die Präparations-/Auslesedeutung. Weder Bosonbanken noch c werden
heimlich aus algebraischer Verfügbarkeit zu experimenteller Kontrolle
erklärt.

